\section{Autres}
\subsection{Exercices exotiques : }
\newpage
\subsubsection{Marche aléatoire isotrope}
Soit $n \in \mathbb{N}$, 

Comme chaque direction peut être choisie de manière équiprobable
il nous suffit de dénombrer le nombre de chemins gagnants qu'on note $G_{2n}$ et on divisera par 
la suite par le nombre de chemin total qu'on note $T_{2n}$. 


En notant un chemin de longueur $2n$ comme une suite dans
$\{B,G,D,H\}$, se donner un chemin gagnant revient à se donner un chemin 
comportant autant de D que de G et autant de B que de H.


Ainsi, choisir un chemin consiste à choisir :
Le nombre de déplacements vers le haut $k \in \llbracket 0,n\rrbracket$

On se convainc qu'un chemin avec $k> n$ déplacements vers le haut
ne peut revenir au centre, car il faudrait au moins $k>n$ déplacements vers le bas, et donc au moins $2k>2n$ déplacements verticaux.

- Quand nous effectuerons les $k$ déplacements vers le haut dans le chemin 
parmi les $2n$ déplacements : $$\binom{2n}{k}$$


- Puis choisir quand nous effectuerons les $k$ déplacements vers le bas parmi les $2n-k$ déplacements restants : 
$$ \binom{2n-k}{k}$$ 

- Il reste $2n-2k$ déplacements à faire, or il faut que le nombre de déplacements à gauche 
soit égal au nombre de déplacements à droite, donc $n-k$ déplacements des deux côtés.

Ainsi il suffit de choisir quand nous nous déplacerons 
à droite et ceux à gauche seront forcés : $$\binom{2n-2k}{n-k}$$

Alors : 

$$ 
\begin{aligned} 
G_{2n} &= \sum_{k=0}^{n}   \binom{2n}{k}\binom{2n-k}{k}\binom{2n-2k}{n-k}
    \\ &= \sum_{k=0}^{n} \frac{(2n)!}{k!(2n-k)!} \frac{(2n-k)!}{k!(2n-2k)!} \frac{(2n-2k)!}{(n-k)!(n-k)!}
    \\ \text{Tout se simplifie bien } &= \sum_{k=0}^{n} \frac{(2n)!}{k!^2(n-k)!^2} = \sum_{k=0}^{n} \frac{n!^2}{k!^2(n-k)!^2} \frac{(2n)!}{n!^2}
    \\ \text{On sort le terme constant qui dérange }&= \frac{(2n)!}{n!^2} \sum_{k=0}^{n} (\frac{n!}{k!(n-k)!})^2 = \frac{(2n)!}{n!^2} \sum_{k=0}^{n} \binom{k}{n}^2
    \\ \text{Par la formule de Vandermonde }&= \frac{(2n)!}{n!^2} \binom{2n}{n}
    \\ \text{Soit finalement : } &= \binom{2n}{n}^2 
\end{aligned}
$$

Par ailleurs, d'après la description précédente, le nombre total de chemin est $$T_{2n}= 4^{2n}$$ 

Donc $$\boxed{P(C_{2n}) = \frac{G_{2n}}{T_{2n}} = \binom{2n}{n}^2 \frac{1}{4^{2n}}}$$

Donnons un équivalent en $+\infty$ : 

$$\begin{aligned}
\displaystyle \text{Par Stirling : } \boxed{P(C_{2n}) = \binom{2n}{n}^2 \frac{1}{4^{2n}} \sim \frac{(\sqrt{4\pi n})(\frac{2n}{e})^{2n}}{(\sqrt{2\pi n}(\frac{n}{e})^n)^44^{2n}}= \frac{1}{n\pi}} \hspace{15px} \textit{(trust les calculs)}
\end{aligned}$$